
\chapter{Diseño e implementación} % Main chapter title

\label{Chapter3}

En este capítulo se explica el desarrollo de la solución alcanzada. Se describen los datos que se procesaron, las estrategias que se probaron, así como el criterio que se usó para compararlas. También se agregan detalles sobre el diseño de la implementación del bot de trading.

\definecolor{mygreen}{rgb}{0,0.6,0}
\definecolor{mygray}{rgb}{0.5,0.5,0.5}
\definecolor{mymauve}{rgb}{0.58,0,0.82}

%%%%%%%%%%%%%%%%%%%%%%%%%%%%%%%%%%%%%%%%%%%%%%%%%%%%%%%%%%%%%%%%%%%%%%%%%%%%%
% parámetros para configurar el formato del código en los entornos lstlisting
%%%%%%%%%%%%%%%%%%%%%%%%%%%%%%%%%%%%%%%%%%%%%%%%%%%%%%%%%%%%%%%%%%%%%%%%%%%%%
\lstset{ %
  backgroundcolor=\color{white},   % choose the background color; you must add \usepackage{color} or \usepackage{xcolor}
  basicstyle=\footnotesize,        % the size of the fonts that are used for the code
  breakatwhitespace=false,         % sets if automatic breaks should only happen at whitespace
  breaklines=true,                 % sets automatic line breaking
  captionpos=b,                    % sets the caption-position to bottom
  commentstyle=\color{mygreen},    % comment style
  deletekeywords={...},            % if you want to delete keywords from the given language
  %escapeinside={\%*}{*)},          % if you want to add LaTeX within your code
  %extendedchars=true,              % lets you use non-ASCII characters; for 8-bits encodings only, does not work with UTF-8
  %frame=single,	                % adds a frame around the code
  keepspaces=true,                 % keeps spaces in text, useful for keeping indentation of code (possibly needs columns=flexible)
  keywordstyle=\color{blue},       % keyword style
  language=[ANSI]C,                % the language of the code
  %otherkeywords={*,...},           % if you want to add more keywords to the set
  numbers=left,                    % where to put the line-numbers; possible values are (none, left, right)
  numbersep=5pt,                   % how far the line-numbers are from the code
  numberstyle=	iny\color{mygray}, % the style that is used for the line-numbers
  rulecolor=\color{black},         % if not set, the frame-color may be changed on line-breaks within not-black text (e.g. comments (green here))
  showspaces=false,                % show spaces everywhere adding particular underscores; it overrides 'showstringspaces'
  showstringspaces=false,          % underline spaces within strings only
  showtabs=false,                  % show tabs within strings adding particular underscores
  stepnumber=1,                    % the step between two line-numbers. If it's 1, each line will be numbered
  stringstyle=\color{mymauve},     % string literal style
  tabsize=2,	                   % sets default tabsize to 2 spaces
  title=\lstname,                  % show the filename of files included with \lstinputlisting; also try caption instead of title
  morecomment=[s]{/*}{*/}
}

\section{Conjunto de datos y preprocesamiento}

Esta sección describe el conjunto de datos utilizado y las estrategias de preprocesamiento aplicadas para el desarrollo del trabajo. Se detallan las consideraciones sobre la frecuencia de los datos, las técnicas de construcción de barras y el análisis de la importancia de las características.


\subsection{Conjunto de datos}

El trabajo se basó en el análisis de datos del mercado spot de Bitcoin. El mercado spot se define como aquel en el que los activos financieros, en este caso Bitcoin (BTC) frente al Tether (USDT), se compran y venden para su liquidación y entrega inmediata. Las variables disponibles para este mercado incluyen el precio de apertura, el precio de cierre, el valor mínimo y máximo alcanzado, el volumen comercializado en unidades de Bitcoin y en USDT, y la hora exacta en la que se tomó cada medición.

Estos datos se recopilaron desde el 1 de enero de 2020 a través de la plataforma de Binance. Inicialmente, se experimentó con una frecuencia de una hora, la cual presentó la desventaja de acotar el volumen de datos disponibles para el entrenamiento del modelo. Además, se observó que las estrategias probadas con esta frecuencia arrojaban peores resultados. Se infirió que esto se debía, posiblemente, a la menor cantidad de datos disponibles, lo que limitó la capacidad del modelo para aprender patrones robustos. Por lo anterior, se optó por utilizar una frecuencia de un minuto, lo que generó una mejora significativa en el desempeño de las estrategias estudiadas al proporcionar un conjunto de datos más denso y amplio.

\subsection{Estrategias de preprocesamiento}

A continuación, se detalla el preprocesamiento realizado, común a la mayoría de los modelos explorados y basado en las recomendaciones de Marcos López de Prado para finanzas cuantitativas.

  Para asegurar la estacionariedad de las series, se aplica la diferenciación fraccional, una técnica avanzada recomendada en la literatura explorada. Dado que las series de precios financieras no son estacionarias, esta técnica busca un orden mínimo de diferenciación (d), que puede ser un número fraccionario (por ejemplo, d=0.4), para hacer la serie estacionaria (según la prueba de Dickey-Fuller aumentada) sin eliminar por completo la información de la serie original. Esto contrasta con la diferenciación clásica (d=1) que puede borrar demasiada ``memoria''.

Para el preprocesamiento de los datos, se aplicaron diversas estrategias con el fin de optimizar su calidad y relevancia para el modelo:

\begin{itemize}
    \item Conversión a micro Bitcoins (uBTC): se realizó una conversión de BTC a uBTC, una práctica estándar en el análisis de criptomonedas para manejar unidades más pequeñas y facilitar los cálculos.
    \item Agrupación por volumen: se exploraron formas de agrupar los datos no por intervalos de tiempo fijos, sino por volumen en unidades constantes tanto de BTC como de USDT. Adicionalmente, se probaron las volume imbalance bars y las tick imbalance bars. Las volume imbalance bars se construyen cuando se alcanza un volumen neto desequilibrado entre compras y ventas, y las tick imbalance bars se forman con un desequilibrio neto en el número de ticks (cambios de precio). Después de una investigación exhaustiva de la literatura y pruebas empíricas, se decidió continuar con la agrupación de barras por volumen fijo. Esta elección se fundamentó en que las barras de volumen aportan una mejora notable en las propiedades estadísticas de la distribución del precio de Bitcoin, ya que reducen la autocorrelación y la heterocedasticidad, lo que las hace más adecuadas para el entrenamiento de modelos de machine learning. Se adoptó un límite de volumen de 50.000 uBTC, que proporcionó un equilibrio satisfactorio entre la cantidad de puntos de datos obtenidos y la reducción de la varianza correspondiente a cada barra.
\end{itemize}

Finalmente, antes del entrenamiento, se aplican transformaciones estándar como el escalado y la reducción de dimensionalidad. El escalado, específicamente \texttt{StandardScaler}, estandariza las características a una media de 0 y varianza de 1, crucial para modelos sensibles a la escala (como SVM o redes neuronales), lo que asegura  que el escalador se ajuste únicamente en los datos de entrenamiento para evitar fugas de información. Opcionalmente, mediante un hiperparametro,  se utiliza el análisis de componentes principales (PCA) para reducir la dimensionalidad y ortogonalizar las características, lo que puede beneficiar a algunos modelos y mejorar la interpretabilidad.

En síntesis, el preprocesamiento representa un pipeline sofisticado que transforma los datos brutos de precios en un formato estructurado, estacionario y etiquetado de manera robusta, listo para ser utilizado por los modelos de aprendizaje de máquina. La figura \ref{fig:preprocessing_pipeline} ilustra este proceso.

\begin{figure}[H]
  \centering
  \includegraphics[width=\textwidth]{"Figures/preprocessing_pipeline.png"}
  \caption{Diagrama del pipeline de preprocesamiento de datos.}
  \label{fig:preprocessing_pipeline}
\end{figure}

%% \clearpage
\subsection{Ingeniería de características y análisis de importancia}

Una vez estructuradas las barras, se procede con la ingeniería de características, donde se construye un conjunto extenso de atributos para los modelos. Esto incluye indicadores técnicos.

Para la ingeniería de características, se generó una extensa lista de indicadores técnicos clásicos como el RSI, el Oscilador Asombroso (AO), Williams \%R y ROC para momentum; MACD, ADX, Aroon, CCI y Vortex para tendencia; y las Bandas de Bollinger, Canales de Keltner y Canales de Donchian para volatilidad y fluctuación. Adicionalmente, se incorporan características basadas en retornos porcentuales (\texttt{pct\_change}) y sus retardos. En el pipeline específico de Chronos, el modelo se utiliza para generar embeddings a partir de ventanas de la serie de precios. Estos embeddings, que encapsulan patrones complejos, sirven como características adicionales para un modelo tabular, como XGBoost. Adicionalmente, se hizo uso de técnicas de reducción de dimensionalidad como el análisis de componentes principales (PCA). Para determinar las características más relevantes para el modelo final, se realizó un análisis de importancia en el cual se utilizó un conjunto de métricas, entre ellas: la disminución media de impureza (MDI), la disminución media de exactitud (MDA) y la importancia de característica individual (SFI). También se analizaron las funciones de autocorrelación (ACF) y autocorrelación parcial (PACF) de los componentes principales para identificar dependencias temporales.

\section{Estrategias de referencia}

Para establecer un punto de comparación robusto, se tomó como principal referencia el trabajo de Guilherme Palazzo \cite{Palazzo:2023}. Su metodología destaca por la aplicación rigurosa de técnicas avanzadas de machine learning financiero, y varios de sus componentes clave fueron adoptados en este trabajo para construir una base sólida sobre la cual evaluar la arquitectura Transformer propuesta.

Los elementos metodológicos que se incorporaron de dicho trabajo son los siguientes:
\begin{itemize}
    \item Muestreo basado en volumen: se adoptó la técnica de agrupar los datos en barras de volumen constante en lugar de intervalos de tiempo fijos. Este enfoque mitiga problemas comunes en las series de tiempo financieras, como la heterocedasticidad, y mejora la calidad de los datos para el entrenamiento de modelos.
    \item Formulación como clasificación: al igual que en el trabajo de referencia, el problema se formuló como una tarea de clasificación para predecir la dirección del movimiento del precio (alza o baja), en lugar de intentar predecir el valor exacto del precio.
    \item Inspiración en el etiquetado dinámico: la estrategia de etiquetado se inspiró en el enfoque de Palazzo de no utilizar horizontes de tiempo fijos. En este trabajo, dicha idea se implementó a través del método de la barrera triple, que se adapta dinámicamente a la volatilidad del mercado.
    \item Validación de modelos basados en árboles: el éxito de XGBoost en el trabajo de referencia sirvió como validación para su uso como un componente eficaz en la arquitectura final, en conjunto con los embeddings generados por Chronos.
\end{itemize}

Al integrar estos componentes, se asegura que la evaluación del rendimiento del modelo Transformer se realice contra una estrategia de referencia competitiva y bien fundamentada en el estado del arte.

\section{Arquitectura transformer propuesta}

Para la arquitectura principal de este trabajo, se utilizó el modelo fundacional Chronos, cuya conceptualización fue detallada en el capítulo 2. En lugar de utilizar el modelo para pronóstico directo, se aprovechó su capacidad para generar representaciones ricas de las series de tiempo.

El enfoque implementado fue el siguiente:
\begin{enumerate}
    \item Se tomaron las series de tiempo de los indicadores técnicos seleccionados.
    \item Cada serie se procesó con un modelo Chronos pre entrenado para extraer los embeddings de sus capas ocultas. Estos embeddings son vectores de alta dimensionalidad que encapsulan la información contextual y los patrones temporales aprendidos por el Transformer.
    \item Los embeddings resultantes de todos los indicadores se concatenaron para formar un único vector de características.
    \item Este vector de características se utilizó como entrada para un modelo tabular más tradicional y computacionalmente eficiente, en este caso, XGBoost o LightGBM, que fue el encargado de realizar la tarea de clasificación final (predecir si el precio subirá o bajará).
\end{enumerate}

Esta estrategia de dos pasos permite combinar el poder de los modelos fundacionales para la extracción de características complejas con la eficiencia y robustez de los modelos basados en árboles para la clasificación en datos tabulares. Finalmente, se experimentó con diferentes tamaños del modelo Chronos (small, base y large) para analizar el compromiso entre el rendimiento predictivo y el costo computacional.

\section{Estrategia de etiquetado}

Para definir el objetivo de predicción del modelo, se evaluaron varias técnicas de etiquetado. Tras analizar los métodos de horizonte fijo y etiquetado dinámico, se adoptó el método de la barrera triple (Triple-Barrier Method). Esta técnica, recomendada en la literatura para mercados financieros, ofrece un enfoque más adaptativo y robusto al definir las etiquetas en función de la volatilidad del activo, con barreras para ganancias, pérdidas y tiempo de expiración.
Para etiquetar un punto se puede elegir una etiqueta independiente para la barrera vertical o agruparla con las barreras horizontales según si el punto final es mayor o menor que el inicial, si es menor se agrupa con la etiqueta correspondiente a la barrera inferior y, si es mayor, con la superior. Esta alternativa permite mantener las clases a predecir balanceadas con lo cual fue la seleccionada.

En la figura \ref{fig:triple_barrier}  se ve cómo a partir de un punto de Incepción (inicio de la operación), se definen tres límites: una barrera de ganancia superior (púrpura) y una barrera de pérdida inferior (cian) que determinan si la operación se clasifica como positiva o negativa si el precio las toca, y una barrera de tiempo vertical (negra) que sirve de vencimiento, etiqueta la posición por su rendimiento hasta ese momento si ninguna barrera horizontal fue alcanzada previamente.

\begin{figure}[H]
  \centering
  \includegraphics[width=	\textwidth]{"Figures/triple-barrier.png"}
  \caption{Aplicación del método de etiquetado de triple barrera sobre un gráfico de precios.\protect}
  \label{fig:triple_barrier}
\end{figure}
%% \footnotetext{}

\section{Framework de backtesting y evaluación}

Para asegurar una evaluación rigurosa y libre de sesgos, se implementó un motor de backtesting a medida en Python. Esta decisión se tomó debido a la necesidad de utilizar metodologías avanzadas como la validación cruzada combinatoria con purga y embargo (\textit{Combinatorially Purged Cross-Validation}). Este método es esencial para prevenir la fuga de información y el sobreajuste al evaluar modelos en series de tiempo financieras, lo que garantiza que el rendimiento medido sea una estimación realista del comportamiento del modelo en datos futuros.

\subsection{Optimización de hiperparámetros: búsqueda bayesiana con Optuna}

La optimización de hiperparámetros es una de las etapas más críticas en la construcción de un modelo de machine learning robusto, especialmente en el dominio financiero donde el riesgo de sobreajuste es muy alto. El enfoque adoptado en este proyecto se centró en la eficiencia y la rigurosidad estadística, y separa estrictamente la fase de optimización de la evaluación final.

Para evitar esto, la búsqueda de hiperparámetros se realizó  únicamente con el conjunto de entrenamiento, y se validó el rendimiento de cada configuración mediante purged k-fold cross-validation. El conjunto de prueba final (hold-out set) se mantuvo completamente aislado y solo se utilizó una única vez al final de todo el proceso para evaluar el modelo definitivo.

Al tratarse de un problema de clasificación, y por el desbalance de clases inherente a las señales de trading, se utilizó la métrica F1 ponderada por clase para evaluar el desempeño. Esta métrica proporciona una medida balanceada entre la precisión y la exhaustividad del modelo, ajustada a la frecuencia de cada clase, lo que la hace más representativa que la exactitud simple.

Para llevar a cabo la búsqueda de manera eficiente, se utilizó la biblioteca Optuna \cite{optuna:url}. En lugar de recurrir a métodos de fuerza bruta como la búsqueda en grilla (Grid Search) o a la ineficiente búsqueda aleatoria (Random Search), Optuna implementa un enfoque de optimización Bayesiana.

Cada prueba individual dentro del estudio de Optuna consistió en un proceso completo y computacionalmente intensivo:

\begin{enumerate}
   \item Sugerencia de hiperparámetros: Optuna propone una nueva combinación de hiperparámetros (ej. \texttt{volume\_threshold}, \texttt{tau}, \texttt{n\_estimators}, etc.).
   \item Ejecución del pipeline completo: con estos hiperparámetros, se ejecuta el pipeline desde cero sobre el conjunto de entrenamiento:
     \begin{enumerate}
       \item Se estructuran los datos en barras de volumen.
       \item Se generan las características (features).
       \item Se crean las etiquetas (labels) según el método definido.
       \item Se ejecuta el proceso completo de purged k-fold cross-validation de n folds.
     \end{enumerate}
   \item Cálculo de la métrica objetivo: se calcula el F1-Score ponderado para cada una de las validaciones cruzadas.
\end{enumerate}

La figura \ref{fig:hyperparameter_pipeline} ilustra este proceso iterativo.

\begin{figure}[ht]
  \centering
  \includegraphics[width=0.8\textwidth]{"Figures/hyperparameter_pipeline.png"}
  \caption{Diagrama del proceso de optimización de hiperparámetros con Optuna.}
  \label{fig:hyperparameter_pipeline}
\end{figure}


Este procedimiento asegura que la evaluación de cada conjunto de hiperparámetros fuera robusta y no dependiera de una única partición de datos.

Una vez que Optuna finaliza todas las pruebas, se identifica el ``mejor trial'', es decir, la combinación de hiperparámetros que produjo el F1-Score promedio más alto. Estos hiperparámetros se consideraron óptimos.

Finalmente, se entrenó un único modelo con estos hiperparámetros óptimos sobre la totalidad del conjunto de entrenamiento. Este modelo final y definitivo fue el que se evaluó, por única vez, en el conjunto de prueba (hold-out) para obtener una estimación realista y sin sesgos de su rendimiento esperado en datos no vistos.

La evaluación de los modelos se centró exclusivamente en su rendimiento en el backtest, lo que separa estrictamente esta fase de la optimización de hiperparámetros. Dicha búsqueda se realizó una única vez durante la etapa de investigación y validación del modelo, con un conjunto de datos distinto al de la evaluación final.


\section{Modelos explorados}

Durante la fase de investigación, se exploró un amplio espectro de modelos para establecer un punto de referencia y comprender la complejidad del problema.

Inicialmente, se evaluaron estrategias de trading clásicas basadas en indicadores técnicos, como el cruce de medias móviles, las bandas de Bollinger, el MACD y las divergencias en el RSI. Si bien estos métodos son populares, demostraron ser demasiado rígidos, poco adaptativos, y propensos al sobreajuste.

Posteriormente, se probaron modelos estadísticos tradicionales para series temporales, incluyendo ARIMA, SARIMA, y variantes más complejas como ARIMAX con GARCH como variable exógena y Prophet de Facebook. Estos modelos, aunque más sofisticados, lucharon por capturar las dinámicas no lineales y no estacionarias del precio de Bitcoin y no superaron a un desempeño aleatorio.

Luego, se exploraron modelos que sean \texttt{drop in replacements} al marco presentado por \cite{Palazzo:2023}. Para seguir con la recomendación de López de Prado sobre la idoneidad de los modelos basados en árboles tipo random forest para datos financieros. Además, se exploraron alternativas a XGBoost. Se utilizó la biblioteca AutoGluon, que permite entrenar y comparar de forma automática un conjunto de modelos tabulares de última generación, incluyendo diferentes configuraciones de redes neuronales y ensambles, para encontrar un posible reemplazo al modelo de Palazzo.

Finalmente, en cuanto a la arquitectura principal de Chronos, se experimentó con diferentes tamaños del modelo T5 (small, base y large) para analizar el compromiso entre el rendimiento predictivo y el costo computacional, se busca el equilibrio óptimo para una implementación práctica.

\section{Diseño del bot}
\label{subsec:bot-design}
El bot de trading se diseñó con una arquitectura modular para facilitar su configuración, mantenimiento y extensibilidad. Se divide en cuatro componentes principales que operan de forma desacoplada:

\begin{enumerate}
    \item Orquestador: es el punto de entrada de la aplicación. Su función es instanciar, configurar y conectar los demás componentes. Inicia el ciclo de operación del bot.
    \item Motor del bot: contiene el bucle principal de ejecución. En cada iteración, orquesta el flujo de trabajo: solicita los datos más recientes del mercado al módulo de intercambio, pasa estos datos al módulo de estrategia para obtener una señal de trading y, si la señal lo indica, instruye al módulo de intercambio para que ejecute la orden correspondiente.
    \item Módulo de exchange: abstrae toda la comunicación con la API del exchange de criptomonedas. Es responsable de obtener datos de mercado (precios, volúmenes) y de ejecutar órdenes (compra, venta). Su diseño intercambiable permite adaptar el bot a diferentes plataformas (por ejemplo, Binance, KuCoin) con mínimas modificaciones.
    \item Módulo de estrategia: carga el modelo de machine learning pre-entrenado (en este caso, el ensamble de Chronos y XGBoost) y lo se utiliza para procesar los datos de mercado y generar una señal de trading: comprar, vender o mantener.
\end{enumerate}

En esencia, el bot opera en un ciclo continuo: recolecta datos, predice una señal con el modelo y actúa sobre esa predicción. Esta separación de responsabilidades garantiza que la lógica de la estrategia de trading esté completamente aislada de los detalles de implementación de la comunicación con el mercado.
de la comunicación con el mercado.
A continuación se proporciona una descripción y vista técnica y exhaustiva de cómo interactúan los cuatro componentes desacoplados (orquestador, motor, exchange y estrategia) a lo largo del ciclo de vida de la aplicación.

\subsubsection{Fase de inicio y orquestación}

Esta fase inicializa el sistema y demuestra el rol del orquestador como punto de entrada de la aplicación, encargándose de instanciar y conectar las piezas. Véase la figura \ref{fig:bot-seq-startup1} y \ref{fig:bot-seq-startup2}.

\begin{figure}[H]
  \centering
  \includegraphics[width=1.0\textwidth]{"Figures/seq_1_startup.png"}
  \caption{Diagrama del proceso de inicialización del bot.}
  \label{fig:bot-seq-startup1}
\end{figure}
\begin{figure}[H]
  \centering
  \includegraphics[width=1.0\textwidth]{"Figures/seq_1b_startup.png"}
  \caption{Diagrama del proceso de inicialización del bot.}
  \label{fig:bot-seq-startup2}
\end{figure}

A continuación se describen los pasos de la fase de inicialización:

\begin{itemize}
\item Con la interacción inicial del usuario: el proceso comienza cuando el usuario introduce el comando bot \texttt{start --port 5000} a través de la interfaz de línea de comandos (cli), alojada en main.py.\\

\item  Creación del servidor en segundo plano: la cli no ejecuta el bot de forma bloqueante. En su lugar, delega la tarea al módulo subprocess para levantar un servidor (\texttt{run\_server}) en el puerto 5000.\\

\item  Confirmación de inicio: la cli devuelve inmediatamente un mensaje al usuario y confirma que el servidor ha iniciado y proporciona el identificador del proceso (pid). Así libera la terminal.\\

\item  Inicialización de servicios concurrentes: dentro del subproceso del servidor, se activan dos componentes críticos de manera paralela. En primer lugar se levanta una aplicación flask (app.run) que servirá para la comunicación y monitoreo externo. En segundo, se inicia un hilo dedicado al bot (botthread) mediante \texttt{thread(start\_bot\_thread).start()}, con esto separa la ejecución web de la lógica de trading.\\

\item Carga de configuraciones: el botthread invoca a la función encargada de leer los parámetros operativos \texttt{load\_config(``config.yaml'')},  y de devolver un diccionario de configuración.\\

\item  Instanciación del módulo de exchange: con su diseño intercambiable, el bot evalúa la configuración recién cargada. Si existe un archivo mock (\texttt{mock\_file}), se instancia un cliente de prueba (mockbinanceclient) para simulaciones. De lo contrario, asume un entorno de producción o testnet y crea un cliente real de binance (binanceclient) con las claves api proporcionadas.\\

\item  Instanciación del módulo de estrategia: el hilo delega la creación del modelo analítico a una fábrica de estrategias (strategyfactory). Dependiendo del nombre en la configuración, la fábrica instancia la clase correspondiente, y carga los parámetros y el modelo subyacente.\\

\item  Arranque del motor del bot: finalmente, el botthread inyecta la configuración, el cliente de exchange y la estrategia recién creados dentro de la instancia principal del bot (el motor) y ejecuta su método principal: bot.run().\\

 \end{itemize}


\subsubsection{Fase de ejecución continua}

Esta es la fase central operada por el motor del bot, caracterizada por un bucle principal que se repite periódicamente cada tick\_interval segundos mientras el flag bot.running sea verdadera. Refleja el ciclo de ``recolectar, predecir y actuar'' ilustrado en la figura \ref{fig:bot-seq-loop}.

\begin{figure}[H]
  \centering
  \includegraphics[width=1.0\textwidth]{"Figures/seq_2_trading_loop.png"}
  \caption{Diagrama del bucle de principal del bot de trading.}
  \label{fig:bot-seq-loop}
\end{figure}

Como vemos en la figura, la ejecución continua consiste en los siguientes pasos:

\begin{itemize}
\item  Recolección de datos: en cada iteración, el bot solicita al exchange los datos históricos del mercado con una invocación a \texttt{get\_historical\_klines} con parámetros de símbolo, marco temporal límite (timeframe). El exchange devuelve una lista de velas (klines).

\item  Procesamiento de datos: el bot toma la lista cruda de velas y construye un dataframe, asegurándose de castear correctamente los tipos de datos para que puedan ser consumidos por el modelo matemático.

\item
  Gestión de riesgos (stop-loss): antes de buscar nuevas señales, el sistema verifica si existe una posición abierta en el mercado. Si es así, el bot calcula internamente el porcentaje de pérdidas y ganancias (pnl).

\item Control de riesgo: si el rendimiento es inferior al umbral de riesgo permitido, el bot ejecuta inmediatamente una orden de venta (\texttt{execute\_sell()}), registra que se activó el stop-loss y finaliza la iteración actual. Después salta directamente al siguiente tick sin consultar a la estrategia.

\item  Predicción de señal: si el stop-loss no se activa, el bot pasa el dataframe procesado al módulo de estrategia mediante la función \texttt{get\_signal(df)}. Tras evaluar los datos, la estrategia devuelve un veredicto direccional estricto: ``buy'' (comprar), ``sell'' (vender) o ``hold'' (mantener).

\item  Ejecución de órdenes de mercado: el bot evalúa la señal de la estrategia cruzada con su estado actual:

\item  Flujo de compra: si la señal es ``buy'' y no hay ninguna posición abierta, el bot le consulta a la estrategia el tamaño de la orden (\texttt{get\_order\_size()}). Tras recibir la cantidad, se ejecuta la orden de compra \\(\texttt{execute\_buy(latest\_price)}) a través del módulo de exchange y actualiza sus balances y estado de posición internamente.

\item  Flujo de venta: si la señal es ``sell'' y el bot ya posee una posición abierta, procede a ejecutar la venta (\texttt{execute\_sell()}) e igualmente actualiza el balance de la cuenta y su estado de posición.

\item  Pausa táctica: al finalizar la lógica operativa, el bot suspende su ejecución mediante \texttt{sleep(tick\_interval)} para evitar saturar las apis del exchange y esperar a la formación de nuevos datos de mercado antes de reiniciar el bucle.

\end{itemize}

\subsubsection{Fase de monitoreo remoto}

El diseño, como se puede ver en la figura \ref{fig:bot-seq-monitoring}, permite auditar el estado del bot en tiempo real sin interrumpir el bucle de trading, gracias a la aplicación web paralela.

\begin{figure}[H]
  \centering
  \includegraphics[width=1.0\textwidth]{"Figures/seq_3_monitoring.png"}
  \caption{Diagrama del proceso de monitoreo del bot de trading.}
  \label{fig:bot-seq-monitoring}
\end{figure}

A continuación se describen los pasos del proceso de monitoreo:

\begin{itemize}

\item  Petición de estado: el usuario ejecuta bot \texttt{status --port 5000} en la cli.

\item  Consulta http: la cli transforma este comando en una petición get dirigida a la ruta /status expuesta por la aplicación flask.

\item  Lectura interna: el servidor flask interactúa con el espacio de memoria del bot, leyendo variables críticas de estado en tiempo real como bot.running (si está activo), la posición actual en el mercado y el balance de la cuenta.

\item  Respuesta al usuario: el bot transfiere este diccionario de estado a flask, que a su vez lo serializa en formato json (conteniendo estatus, posición y balance) y lo envía de regreso a la cli. Finalmente, la cli imprime esta información de forma legible para el usuario.

\end{itemize}

\subsubsection{Fase de apagado}

El sistema contempla mecanismos robustos para finalizar el proceso del bot de manera segura. El diagrama \ref{fig:bot-seq-shutdown} expone dos vías para detener la operación:
\begin{figure}[H]
  \centering
  \includegraphics[width=1.0\textwidth]{"Figures/seq_4_shutdown.png"}
  \caption{Diagrama del proceso de apagado del bot de trading.}
  \label{fig:bot-seq-shutdown}
\end{figure}

A continuación se describen los pasos del proceso de apagado:

\begin{itemize}
\item  Apagado forzado por cli: el usuario emite el comando bot stop. La cli lee el archivo que contiene el identificador del proceso (bot.pid) e interactúa directamente con el sistema operativo mediante el envió de una señal de terminación (\texttt{os.kill(pid, sigterm)}) al servidor para que aborte la ejecución.

\item  Apagado elegante vía api: como método alternativo anotado en el diagrama, se puede enviar una petición post a la ruta /shutdown a través de flask. En este escenario, flask se comunica con el bot con una llamada al método bot.stop(), el cual cambia internamente la variable self.running a false. Esto permite que el bucle principal finalice su iteración actual de manera limpia, con la lubricación de los hilos y el cierre del servidor de forma orquestada y sin corrupción de estado.
\end{itemize}

\section{Diseño del sistema}

El sistema esta estructurado siguiendo la temática de pipeline de investigación propuesta por López de Prado en Advances in financial machine learning:
\begin{enumerate}
\item \texttt{data\_analysis}: este directorio implementa el paso de curación de datos y de análisis de característica, aquí se proveen rutinas para recopilar, limpiar, indexar, almacenar, y ajustar los datos para entregarlos en un formato útil para el resto de la cadena de producción y también para transformar datos brutos en señales informativas con valor predictivo. En particular se implementan, en el paquete \texttt{bar\_aggregation},  funciones para agregar barras de tiempo en barras de volumen. Estas incluyen estrategias por volumen de USDT o de BTC. También se encuentran otras utilidades para el preprocesamiento de datos, en el subpaquete \texttt{data\_analysis}, como ser la diferenciación fraccionaria, operadores de medias móviles. Aquí se encuentran los mas de 25 indicadores de trading usados en este trabajo, como ser momentum, relative strenght index, MACD, entre otros. Las funciones para el etiquetado, descriptas anteriormente, como son, por horizonte fijo, por horizonte ajustado por volatilidad, y de barrera triple, se hacen disponible en el subpaquete \texttt{labeling}.

\item \texttt{modeling}: en este paquete viven los algoritmos de inversión que utilizan los datos provistos por las subrutinas encontradas en el paquete de \texttt{data\_analysis} así también como las subrutinas para procesar dichos datos y transformarlos en variables independientes para los modelos.  Para los propósitos de este trabajo aquí es donde se encuentran las cuatro categorías de estrategias que se desarrollaron: modelos basados en indicadores de trading, modelos basados en modelos estadísticos tradicionales, modelos basados en predictores clásicos de machine learning, y modelos basados en deep learning, en particular, utiliza la arquitectura de transformers.
  Para su evaluación aquí también se implementa el algoritmo de k-fold purged cross validation  ya descripto anteriormente. En adición a esto también se encuentran utilidades para la reutilización de código dedicado a la interacción con la biblioteca mlflow dedicado a guardar los resultados obtenidos en los experimentos realizados.
  Con el objetivo de coordinar las interacciones previamente mencionadas y de estandarizar como se alimentan los algoritmos con datos, como se procesan y se orquestaran para su evaluación. Centralizar y reutilizar esta lógica asegura que siempre se van a tomar las mismas medidas y con el mismo criterio para todos los algoritmos estudiados. Va a evitar que, al reescribir código similar, se introduzcan variantes o errores conceptuales a la hora de separar los datos de entrenamiento y de prueba. Como contrapartida, esto obliga a, sin incurrir en una complejidad excesiva, destilar  una interfaz lo suficientemente flexible como para ser utilizada por diversas categorías de modelos con distintos requerimientos en lo que respecta a los dominios de estructurado de datos, etiquetado, e ingeniería de características.
\item \texttt{backtesting}:  aquí se evalúa la rentabilidad de la estrategia bajo diversos escenarios simulados y estiman la probabilidad de sobreajuste del backtest teniendo en cuenta el número de ensayos y errores realizados. Los principales subpaquetes a nombrar son \texttt{cpcv} y \texttt{cpcv\_runner}. En el primero se implementan las funciones puras que calculan los splits combinatorios, la lógica de embargo, y la reconstrucción de caminos para el algoritmo de combinatorially purged cross validation. En el segundo se implementan funciones que unifican la lógica par aplicar CPCV a los modelos existentes.

\item \texttt{app}: en este paquete vive el bot de trading tal cual descripto en la sección \ref{subsec:bot-design}.
\end{enumerate}


\section{Pipeline de modelos}
\label{subsec:ml-pipeline}
La abstracción central utilizada en el desarrollo de las pruebas para la optimización fue el concepto de pipeline. Como ya se dijo, este fue con el objetivo de centralizar la coordinación de las áreas principales que componen la investigación de estrategias de inversión según López de Prado. En la figura \ref{fig:ml-pipeline-flow} se ve conceptualmente el flujo de datos pasar a través del pipeline.
Describe dos fases. La fase inicial transforma datos crudos en un formato tabular. El proceso comienza con el estructurado de datos, que genera barras de información. Estas barras alimentan dos procesos paralelos: la ingeniería de características, que crea variables predictivas, y el etiquetado y ponderación, que define las variables objetivo. Luego, ambas ramas convergen en el paso de alineamiento: lo que garantiza que todas las matrices compartan los mismos índices temporales para formar el dataset final.

La segunda fase utiliza este dataset para modelar y validar la estrategia. Los datos ingresan al módulo para dividir folds: aplica una validación cruzada purgada para evitar filtraciones de información. Este proceso inicia un ciclo repetitivo de n folds. Dentro de cada iteración secuencial ocurren cuatro pasos: el escalado ajusta los datos, la reducción de dimensiones extrae la información relevante, el entrenamiento y predicción ajusta el algoritmo para generar pronósticos, y finalmente la evaluación mide el rendimiento frente a datos reales. Así se conecta el procesamiento analítico con una validación rigurosa del modelo.

\begin{figure}[H]
  \centering
  \includegraphics[width=1.0 \textwidth]{"Figures/ml-pipeline-flow.drawio.png"}
  \caption{Flujo de datos en distintas etapas de procesamiento del pipeline.}
  \label{fig:ml-pipeline-flow}
\end{figure}

Con esta estructura de base se pudieron evaluar los distintos modelos y estrategias de forma consistente mediante el uso de herencia como se ve en la figura \ref{fig:ml-pipeline-flow-abstract}. 
\begin{figure}[H]
  \centering
  \includegraphics[width=1.0\textwidth]{"Figures/ml-pipeline-flow-abstract-mlpipeline.drawio.png"}
  \caption{Estructura jerárquica de herencia de clase para AbstractMLPipeline.}
  \label{fig:ml-pipeline-flow-abstract}
\end{figure}


\section{Estructura de pruebas realizadas con Chronos}

Se aprovecha la flexibilidad brindada por la estructura de clases descripta en la sección \ref{subsec:ml-pipeline} (véase figura \ref{fig:ml-pipeline-chronos}) para realizar una serie de pruebas centradas en la familia de modelos Chronos. Las pruebas se pueden dividir conceptualmente en dos grandes categorías: las que usan Chronos para realizar la predicción y las que usan chronos para hacer transfer learning.

\begin{figure}[H]
  \centering
  \includegraphics[width=1.0 \textwidth]{"Figures/ml-pipeline-flow-chronos-pipeline.drawio.png"}
  \caption{Estructura jerárquica de herencia de clase para los experimentos realizados con Chronos.}
  \label{fig:ml-pipeline-chronos}
\end{figure}

En la primera categoría se utiliza el modelo chronos como regresor o como clasificador. En forma de regresor se tiene clases \texttt{PalazzoChronosPipeline} que trata de predecir el precio directamente. Del lado de los clasificadores se tienen \texttt{PalazzoChronosClassificationPipeline} que realiza una regresión sobre el precio de cierre pero luego solo se toma si es mayor o menor al cierre anterior, en efecto transformándolo en un clasificador; y \\\texttt{PalazzoChronosBinaryClassificationPipeline} al que, en vez de alimentarlo con los precios de cierre, se le proveen directamente las etiquetas +1 y -1 que indican si el precio sera mayor o menor que el anterior.

En la categoría de transfer learning se tiene a la familia de los \\ \texttt{ChronosFeaturePipelines} que modifican el pipeline de inspirado en el trabajo de Palazzo, que se usó como línea de base, para agregar los embeddings de chronos como características para que sean usadas por un modelo tabular.

En ambas categorías se realizaron pruebas con inferencia de zero-shot y también luego de hacer fine-tuning. Para éste último se utilizó una estrategia de low rank adaptation lo que hizo factible la tarea en el hardware disponible.

